\appendix
\section{Holomorphic references and boundary interpolation}
\label{app:reference}

This appendix makes explicit two analytic ingredients of the earlier
finite comparisons \cite{note02}. They are separate from the signed
source estimates used in the endpoint program.

\begin{proposition}[Minimum norm for a zero-free square holomorphic map]
\label{a:minimum}
Let $\Omega\Subset\mathbb C^n$ be a bounded domain with piecewise smooth
boundary, and let $F$ be holomorphic near $\overline\Omega$, with
values in $\mathbb C^n$. Suppose
$b=\min_{\partial\Omega}\|F\|>0$. If
$\deg(F,\Omega,0)=0$, then $F$ has no zero in $\overline\Omega$ and
\[
 \min_{\overline\Omega}\|F\|=b.
\]
\end{proposition}
\begin{proof}
For $\|v\|<b$, the fiber $F^{-1}(v)\cap\Omega$ is a compact analytic
subset of the open Euclidean domain: it cannot approach the boundary.
Such a subset is finite. Indeed, on a compact irreducible analytic
component the maximum principle makes every holomorphic coordinate
function constant, excluding positive dimension; compactness then
makes the discrete set finite. At each isolated holomorphic preimage
the local degree is a positive integer, the analytic multiplicity.
The degree at $v$ is their sum. It is constant for $\|v\|<b$, since
that ball misses $F(\partial\Omega)$. Degree zero at $0$ therefore
implies that every such fiber is empty. This proves the claimed
interior lower bound, and compactness of the boundary gives equality.
\end{proof}

\begin{corollary}[Transfer of the same certified margin]
\label{a:transfer}
Under the hypotheses of Proposition~\ref{a:minimum}, let $G$ be another
square holomorphic map near $\overline\Omega$ and suppose
$\sup_{\partial\Omega}\|G-F\|\leq e<b$. Then
\[
 \|G\|\geq b-e\quad\hbox{on }\overline\Omega,
 \qquad \deg(G,\Omega,0)=0.
\]
\end{corollary}
\begin{proof}
Apply the scalar maximum-modulus principle to every unit linear
functional of $G-F$. Taking the supremum over those functionals gives
$\sup_{\overline\Omega}\|G-F\|\leq e$. The norm inequality follows
from Proposition~\ref{a:minimum}; the straight boundary homotopy
$F+s(G-F)$ is nonzero and preserves degree.
\end{proof}

The same supplied reference-degree proof, boundary floor and error
certificate thus yield a direct interior joint margin without additional
source acquisition. This is a comparison of those data, not a lower
bound for the cost of every boundary algorithm. Computing the reference
degree has a cost shared by both uses.

The margin is sharp even for a heat jet requiring both components.
On $|\operatorname{Re}t|,|\operatorname{Im}t|<1/10$ and
$|\operatorname{Re}z|,|\operatorname{Im}z|<3/2$, set
\[
 F=(z^2-1-2t,2z),\qquad
 G=(z^2-1-2t+\eta,2z),\quad 0<\eta<4/5.
\]
Writing $s=|z|^2$ and $k=1+2t$ gives
\[
 \|F\|^2\geq(s-|k|)^2+4s
 =|k|^2+(4-2|k|)s+s^2\geq16/25,
\]
because $\operatorname{Re}k\geq4/5$ and $|k|<2$.
Equality holds at $t=-1/10,z=0$.
The identical argument with $k-\eta$ gives exact floor $4/5-\eta$
for $G$, while $\|G-F\|=\eta$. At time zero, the first component
of $F$ vanishes at $z=\pm1$, and its second component at $z=0$.
Neither component by itself supplies the joint floor.

\subsection{A complete interpolation payment}

Suppose a convex boundary simplex is approximated by affine
interpolation of the jets $H_j=\partial_z^jH$, $j=0,1$.
Let $(t_i,z_i)$ be its vertices, $\lambda_i$ barycentric weights,
and $(t,z)=\sum_i\lambda_i(t_i,z_i)$. Put
\[
 V_t=\sum_i\lambda_i|t_i-t|^2,\qquad
 V_z=\sum_i\lambda_i|z_i-z|^2.
\]
If complete moments $M_k(T,R)$ bound all segment derivatives, Taylor's
formula and cancellation of the linear barycentric term give
\begin{equation}
 \left|\sum_i\lambda_iH_j(t_i,z_i)-H_j(t,z)\right|
 \leq I_j:=\frac12\left(M_{j+4}V_t+
 2M_{j+3}\sqrt{V_tV_z}+M_{j+2}V_z\right).
 \label{a:interpolation}
\end{equation}
To check this, the second directional derivative has terms
$H_{j,tt}=H_{j+4}$, $H_{j,tz}=-H_{j+3}$ and
$H_{j,zz}=H_{j+2}$. Integrating the Taylor remainder contributes
$1/2$, and Cauchy--Schwarz bounds the mixed barycentric sum.
For a real coordinate with range length $h$, its variance is at most
$h^2/4$: this follows by taking expectations in
$(X-a)(b-X)\geq0$ and maximizing
$(\mathbb EX-a)(b-\mathbb EX)$. Summing over the two time and
two space coordinates bounds $V_t,V_z$ explicitly.

If the vertex errors in $H_j$ are at most $\delta_j$, convexity keeps
the interpolated acquisition errors at most $\delta_j$. Thus the
full Euclidean jet error on the simplex is at most
\begin{equation}
 e_\sigma=\sqrt{(\delta_0+I_0)^2+(\delta_1+I_1)^2}.
 \label{a:fullerror}
\end{equation}
A certified lower bound $b_\sigma$ for the norm on the entire affine
image, with $e_\sigma<b_\sigma$, gives a nonzero straight boundary
homotopy. Vertex norms alone are insufficient for $b_\sigma$; for
example, enclosing the image convex hull in a coordinate box permits
a lower bound by that box's distance from zero.

For completeness, a finite piecewise affine degree computation can
cone a coherently triangulated box boundary to its center, assigning
the artificial target value zero to the cone vertex. This is an
extension of the boundary map, not an interior approximation to $G$.
For each four-simplex form the domain and target matrices
$X=[x_0\ x_1\ x_2\ x_3]$ and
$Y=[y_0\ y_1\ y_2\ y_3]$ relative to the cone center. A sufficiently
small regular target $v$ contributes $\operatorname{sgn}(\det Y/\det X)$
exactly when $Y^{-1}v$ has positive coordinates with sum less than one.
Choose $\|v\|<\min_\sigma b_\sigma$ and verify that it is on no
simplex image face and has no preimage in a singular simplex. The sum
is the boundary degree, by additivity and homotopy invariance. Rational
linear algebra can verify these finite conditions exactly. The
artificial center uses no evaluation of the actual heat field. Finally,
holomorphic positivity is what makes zero degree exclude a zero of
$(H,H_z)$; that last conclusion would fail for a general real map.
